\documentclass[12pt]{article}

% ==========================================
% PAQUETES
% ==========================================

\usepackage[spanish]{babel}
\usepackage[utf8]{inputenc}
\usepackage[T1]{fontenc}

\usepackage{geometry}
\usepackage{graphicx}
\usepackage{setspace}
\usepackage{titlesec}
\usepackage{float}
\usepackage{booktabs}
\usepackage{amsmath}
\usepackage{caption}

% ==========================================
% CONFIGURACIÓN
% ==========================================

\geometry{margin=1in}
\onehalfspacing

\begin{document}

% ==========================================
% PORTADA
% ==========================================

\begin{titlepage}
    \centering

    % Logo opcional
    % \includegraphics[width=0.25\textwidth]{logo.png}
    % \vspace{1cm}

    {\scshape\LARGE Matemáticas para Ciencia de Datos \par}

    \vspace{2cm}

    {\Huge\bfseries Actividad de Aprendizaje 2.2 \par}

    \vspace{0.7cm}

    {\Large
    Análisis de un Set de Datos Multivariable\\
    con Distribuciones Multimodales
    \par}

    \vspace{2.5cm}

    \begin{flushleft}
    \large
    \textbf{Alumno:} Emilio Izquierdo Montero \\[0.3cm]

    \textbf{Profesor:} Adán Mendoza Soria \\[0.3cm]

    \textbf{Grupo / Semestre:} 8 \\[0.3cm]

    \textbf{Fecha:} 24 de mayo de 2026
    \end{flushleft}

    \vfill

\end{titlepage}


\tableofcontents
\newpage


\section{Introducción}

En esta actividad se realizó el análisis de un conjunto de datos multivariable
utilizando herramientas estadísticas y visuales en Python.

El objetivo principal fue estudiar el comportamiento de distribuciones
multimodales, así como comprender el impacto que tienen sobre las medidas de
tendencia central y el análisis exploratorio de datos.

Para ello se simuló un sistema de monitoreo de servidores con variables como:

\begin{itemize}
    \item Tiempo de respuesta del sistema
    \item Uso de CPU
    \item Consumo de memoria
\end{itemize}

El dataset fue construido utilizando combinaciones de distribuciones normales
para generar comportamientos multimodales y analizar posibles correlaciones
entre variables.

\section{Generación del Dataset}

El dataset fue generado artificialmente mediante Python utilizando las
bibliotecas:

\begin{itemize}
    \item NumPy
    \item Pandas
    \item Matplotlib
    \item Seaborn
\end{itemize}

Se generaron 1000 registros correspondientes al comportamiento de un sistema
de servidores.

La variable principal, correspondiente al tiempo de respuesta, fue creada
mediante la combinación de dos distribuciones normales:

\[
X = 0.6 \cdot N(120,15) + 0.4 \cdot N(300,40)
\]

Esto permitió producir una distribución bimodal, representando dos estados
distintos del sistema:

\begin{itemize}
    \item Servidores en operación normal
    \item Servidores bajo alta carga
\end{itemize}


\section{Medidas de Tendencia Central}

\subsection{Media}

La media representa el promedio de todos los valores de una variable.

\[
\bar{x} = \frac{1}{n}\sum_{i=1}^{n}x_i
\]

La media permite obtener una idea general del comportamiento promedio del
sistema.

\subsection{Mediana}

La mediana corresponde al valor central de un conjunto de datos ordenado.

Es especialmente útil cuando existen valores extremos o distribuciones
asimétricas, ya que es menos sensible a valores atípicos que la media.

\subsection{Moda}

La moda representa el valor o conjunto de valores con mayor frecuencia dentro
de una distribución.

En distribuciones multimodales pueden existir múltiples modas, lo cual indica
la presencia de distintos grupos o comportamientos dentro de los datos.


\section{Distribuciones Multimodales}

Las distribuciones multimodales contienen más de un pico principal dentro de
su histograma.

En este caso, la variable de tiempo de respuesta presentó dos agrupamientos
claramente diferenciados:

\begin{itemize}
    \item Un grupo de baja latencia
    \item Un grupo de alta latencia
\end{itemize}

Esto evidencia que el sistema presenta distintos estados operativos.

La media puede resultar engañosa en estos escenarios, ya que puede ubicarse en
una región donde realmente existen pocos datos.


\section{Resultados Estadísticos}

La Tabla \ref{tab:estadisticas} muestra las principales medidas estadísticas
calculadas para cada variable del dataset.

\begin{table}[H]
\centering

\caption{Medidas estadísticas del dataset}
\label{tab:estadisticas}

\begin{tabular}{lccccc}
\toprule
\textbf{Variable} & \textbf{Media} & \textbf{Mediana} & \textbf{Moda} & \textbf{Varianza} & \textbf{Desv. Est.} \\
\midrule

Tiempo de Respuesta & 191.52 & 152.84 & 117.23 & 7741.34 & 87.98 \\

Uso de CPU & 69.41 & 66.28 & 58.17 & 235.76 & 15.35 \\

Memoria & 1452.61 & 1398.74 & 1204.55 & 132544.82 & 364.06 \\

\bottomrule
\end{tabular}

\end{table}

Los resultados muestran una dispersión considerable en las variables,
especialmente en el tiempo de respuesta y el consumo de memoria.

La diferencia entre media y mediana sugiere la presencia de distribuciones
asimétricas y multimodales.


\section{Histogramas}

Los histogramas permiten visualizar la distribución de cada variable y detectar
la presencia de múltiples agrupamientos dentro de los datos.

\subsection{Histograma del Tiempo de Respuesta}

\begin{figure}[H]
\centering
\includegraphics[width=0.8\textwidth]{/home/emilio/proyectos/adan1/graficas/hist_respuesta.png}
\caption{Histograma del tiempo de respuesta}
\end{figure}

En este histograma se observan claramente dos picos principales, lo que
indica una distribución bimodal.

Esto sugiere la existencia de dos estados operativos distintos en el sistema.


\subsection{Histograma del Uso de CPU}

\begin{figure}[H]
\centering
\includegraphics[width=0.8\textwidth]{/home/emilio/proyectos/adan1/graficas/hist_cpu.png}
\caption{Histograma del uso de CPU}
\end{figure}

El uso de CPU presenta una distribución más dispersa, aunque aún se observan
agrupamientos relacionados con los distintos niveles de carga del sistema.


\subsection{Histograma del Consumo de Memoria}

\begin{figure}[H]
\centering
\includegraphics[width=0.8\textwidth]{/home/emilio/proyectos/adan1/graficas/hist_memoria.png}
\caption{Histograma del consumo de memoria}
\end{figure}

El consumo de memoria presenta una tendencia creciente y una dispersión amplia,
relacionada con el comportamiento del tiempo de respuesta y el uso de CPU.


\section{Gráficos de Dispersión}

\subsection{Tiempo de Respuesta vs Uso de CPU}

\begin{figure}[H]
\centering
\includegraphics[width=0.8\textwidth]{/home/emilio/proyectos/adan1/graficas/scatter_respuesta_cpu.png}
\caption{Tiempo de respuesta vs uso de CPU}
\end{figure}

Se observa una correlación positiva entre ambas variables:
a mayor tiempo de respuesta, mayor uso de CPU.

También se distinguen agrupamientos que podrían representar diferentes estados
de carga del sistema.

% -----------------------------------------------------

\subsection{Tiempo de Respuesta vs Consumo de Memoria}

\begin{figure}[H]
\centering
\includegraphics[width=0.8\textwidth]{/home/emilio/proyectos/adan1/graficas/scatter_respuesta_memoria.png}
\caption{Tiempo de respuesta vs consumo de memoria}
\end{figure}

Existe una relación creciente entre el tiempo de respuesta y el consumo de
memoria.

Los datos muestran la presencia de clusters o conglomerados asociados a
distintos niveles de rendimiento.


\subsection{Uso de CPU vs Consumo de Memoria}

\begin{figure}[H]
\centering
\includegraphics[width=0.8\textwidth]{/home/emilio/proyectos/adan1/graficas/scatter_cpu_memoria.png}
\caption{Uso de CPU vs consumo de memoria}
\end{figure}

La relación entre CPU y memoria también presenta una tendencia positiva,
aunque con mayor dispersión.

Esto indica que ambas variables están parcialmente relacionadas.


\section{Análisis de Resultados}

A partir del análisis estadístico y visual se identificaron varios patrones
importantes dentro del dataset.

Los histogramas permitieron detectar claramente la multimodalidad en el tiempo
de respuesta, evidenciando la existencia de distintos comportamientos dentro
del sistema.

Los gráficos de dispersión mostraron correlaciones positivas entre las
variables analizadas, especialmente entre:

\begin{itemize}
    \item Tiempo de respuesta y uso de CPU
    \item Tiempo de respuesta y consumo de memoria
\end{itemize}

Además, se observaron agrupamientos o clusters que sugieren diferentes estados
operativos de los servidores.

Esto demuestra que las visualizaciones aportan información adicional que no
puede identificarse únicamente mediante medidas numéricas.

Para modelos predictivos, este comportamiento implica que podrían existir
subpoblaciones dentro de los datos, por lo que técnicas como clustering o
segmentación podrían mejorar significativamente el rendimiento de modelos de
machine learning.
\section{Conclusiones}

Los resultados obtenidos demuestran que las medidas estadísticas tradicionales
no siempre son suficientes para describir correctamente un conjunto de datos.

Las distribuciones multimodales pueden ocultar información importante cuando
solamente se utiliza la media como referencia.

Las visualizaciones permitieron detectar patrones, agrupamientos y posibles
correlaciones entre variables, proporcionando una interpretación más completa
del comportamiento del sistema.


\end{document}
